% Generated from common/cover-notes.json by tools/series_cover.py.
\backmatter
\cleardoublepage
\chapter*{About the Cover}
\addcontentsline{toc}{chapter}{About the Cover}
\noindent A regular icosahedron, drawn with one axis through a vertex, one through a face and one through an edge, and with the circles in which its mirror planes meet the view. Euclid constructs the solid in Book XIII of the Elements. Felix Klein's Lectures on the Icosahedron (1884) identified its rotations with the alternating group on five letters, the geometric reason a general quintic has no solution by radicals. The closing chapter of this volume, A Glimpse Beyond: Quaternions and Rotations, is that classification. The drawing was computed by the book's own TikZ cover generator.
